% $Id: implementations.tex 1784 2012-04-27 23:29:31Z  $
% !TEX root = main.tex

\chapter{Verifying CRDT Implementations}\label{chap:implementations}

\section{Introduction}

Once established the foundation (Chapter~\ref{chap:theory}) and the unified framework architecture (Chapter~\ref{chap:framework}), we present different CRDTs implementations that demonstrate the verification capability of our system.

\textbf{Note:} This chapter presents the main theorems for each CRDT. The complete verified implementations, including all lemmas and helper proofs, are available in the repository.\footnote{\url{https://github.com/FLAGlab/Athena-CRDTs/tree/verified-implementation}}

For each CRDT, we demonstrate the following classes of verified properties:

\begin{enumerate}
  \item \textbf{Monotonicity properties:} Operations and merges never decrease over state components. This ensures both local updates and merge operations move upward in the lattice partial order.
    
    \item \textbf{Progress properties:} Updates applied at one replica remain visible after merging with concurrent states. Formally, $\mathit{merge}(s, \mathit{apply}(s, v)) = \mathit{apply}(s, v)$.
    
    \item \textbf{Behavioral guarantees:} CRDT-specific properties like add-wins semantics (OR-Set), deterministic tie-breaking (LWW-Register), or component independence (PN-Counter).
\end{enumerate}

We organize our presentation into three representative CRDTs demonstrating different verification challenges:
\begin{itemize}
  \item \textbf{PN-Counter} (Section~\ref{sec:pncounter}) -- a state-based CRDT with proven properties demonstrating monotonicity and progress guarantees.
  \item \textbf{LWW Register} (Section~\ref{sec:lww}) -- a causal CRDT showing how vector clocks enable deterministic concurrent write resolution.
  \item \textbf{OR-Set} (Section~\ref{sec:orset}) -- a grow-only set with add-wins semantics, demonstrating complex set membership reasoning and re-addition correctness.
\end{itemize}

\textbf{Note}: Complete implementations for all six CRDTs (G-Counter, G-Set, PN-Counter, 2P-Set, OR-Set, and LWW-Register) are available in the code \href{https://github.com/FLAGlab/Athena-CRDTs/tree/verified-implementation}{repository}. We present the most representative here.

\section{PN-Counter: A Showcase of Verified Properties}\label{sec:pncounter}

We focus on the PN-Counter (Positive-Negative Counter) as a representative state-based CRDT. PN-Counter supports both increments and decrements through a pair of GCounters: one for positive updates ($P$) and one for negative updates ($N$). The observed value is computed as $P - N$. 

This design illustrates the state (lattice position) grows monotonically under the merge operation, but the computed value (obtained through subtraction) may decrease. Formal verification is essential to ensure that this separation between state semantics and value computation does not introduce unintended behavior.
 
\subsection{Specification and Properties}

\paragraph{State representation.}
\[
\text{State} = (\text{Ide} \to \mathbb{N}) \times (\text{Ide} \to \mathbb{N}) = P \times N
\]
A pair of finite maps, each associating replicas with natural number counts.

\paragraph{Operations.} The PN-Counter supports the following operations:
\begin{itemize}
  \item \textit{increment}: $\mathit{increment}((P, N), r) = (P[r \mapsto P(r)+1], N)$
  \item \textit{decrement}: $\mathit{decrement}((P, N), r) = (P, N[r \mapsto N(r)+1])$
  \item \textit{merge}: $\mathit{merge}((P_1, N_1), (P_2, N_2)) = (\text{max}_r(P_1(r), P_2(r)), \text{max}_r(N_1(r), N_2(r)))$
  \item \textit{value}: $\mathit{value}((P, N)) = \sum_r P(r) - \sum_r N(r)$
\end{itemize}

The merge operation applies componentwise maximum to both the positive and negative components. The value function performs arithmetic subtraction, which may decrease even though the underlying state (the pair of maps) grows monotonically.

\paragraph{PN-Counter's theorem collection.}

The PN-Counter specification contains \textbf{22 proven theorems}, among which we highlight monotonicity, progress preservation, and semantic distinction:

\begin{table}[h]
  \centering
  \small
  \begin{tabular}{lp{5cm}}
    \toprule
    \textbf{Category} & \textbf{Representative Theorems}\\
    \midrule
    Structural & \texttt{pnc-merge-is-pointwise-max}, \texttt{merge-commutative}\\
    \midrule
    Component Monotonicity & \texttt{positive-component-monotonic}, \texttt{negative-component-monotonic}\\
    \midrule
    Update Monotonicity & \texttt{increment-pos-monotonic}, \texttt{increment-neg-monotonic}\\
    \midrule
    Progress & \texttt{increment-increases-P}, \texttt{decrement-increases-N}\\
    \bottomrule
  \end{tabular}
  \caption{Representative theorems from PN-Counter's verified properties. Progress properties (\texttt{increment-increases-P}, \texttt{decrement-increases-N}) demonstrate that updates increase their respective components, while merge preservation properties (\texttt{merge-preserves-pos-update}, \texttt{merge-preserves-neg-update}) ensure updates survive merging.}
  \label{tab:pncounter-theorems}
\end{table}

These theorems demonstrate different aspects of CRDT correctness: (1) \emph{component monotonicity} ensures merge operations never decrease state components and (2) \emph{progress preservation} ensures local updates are preserved after merging with concurrent states.

First, Listing~\ref{lst:pncounter-monotonic} establishes component monotonicity:

\begin{lstlisting}[language=Athena,style=athena,label={lst:pncounter-monotonic},caption={PN-Counter: Component monotonicity}][b!]
conclude positive-component-monotonic :=
  (forall s1 s2 . (N.<= (first s1) (first (pnc-merge s1 s2))))
  pick-any s1:PNCState s2:PNCState
    (!two-cases
      assume case1 := ((first s2) < (first s1))
        conclude ((first s1) N.<= (first (pnc-merge s1 s2)))
          let {eq1 := (!chain [(first (pnc-merge s1 s2))
                       = (N.Max (first s1) (first s2)) [pnc-merge-is-pointwise-max]
                       = (first s1) [N.Max.less2]]);
               refl := (!uspec N.Less=.reflexive (first s1))}
            (!rcong refl ((first s1) N.<= (first (pnc-merge s1 s2))))
      assume case2 := (~ ((first s2) < (first s1)))
        conclude ((first s1) N.<= (first (pnc-merge s1 s2)))
          (!chain-> [case2
                   ==> ((first s1) N.<= (first s2)) [N.Less=.trichotomy2]
                   ==> ((first s1) N.<= (N.Max (first s1) (first s2)))
                   ==> ((first s1) N.<= (first (pnc-merge s1 s2)))]))
\end{lstlisting}

This proves the positive component never decreases through merge.

Second, Listing~\ref{lst:pncounter-progress} demonstrates progress preservation:

\begin{lstlisting}[language=Athena,style=athena,label={lst:pncounter-progress},caption={PN-Counter: Progress preservation (updates survive merging)}]
conclude merge-preserves-pos-update :=
  (forall s s2 d .
    (first (pnc-apply-update s (inc-pos d))) N.<= 
    (first (pnc-merge (pnc-apply-update s (inc-pos d)) s2)))
  pick-any s:PNCState s2:PNCState d:N
    let {s-updated := (pnc-apply-update s (inc-pos d))}
      (!uspec (!uspec positive-component-monotonic s-updated) s2)
\end{lstlisting}

This establishes that local updates survive merge operations. The proof instantiates the component monotonicity theorem twice: first specializing it to the updated state \texttt{s-updated}, then to the arbitrary state \texttt{s2}. If the positive component never decreases, then an increment to that component will remain visible after any subsequent merge.

Third, Listing~\ref{lst:pncounter-semantic} formalizes semantic distinction:

\begin{lstlisting}[language=Athena,style=athena,label={lst:pncounter-semantic},caption={PN-Counter: State grows but value may decrease}]
conclude state-grows-value-may-decrease :=
  (forall s1 s2 . 
    ((N.<= (first s1) (first s2)) ==> 
     (N.<= (second s1) (second s2)) ==>
     (pnc-merge s1 s2) = s2))
  pick-any s1:PNCState s2:PNCState
    assume h1 := ((first s1) N.<= (first s2))
      assume h2 := ((second s1) N.<= (second s2))
        let {max-f := (!two-cases ...);
             max-s := (!two-cases ...)}
          (!chain [(pnc-merge s1 s2)
                 = (pr (N.Max (first s1) (first s2))
                       (N.Max (second s1) (second s2)))
                 = (pr (first s2) (second s2)) [max-f max-s]
                 = s2])
\end{lstlisting}

Together, these three theorems (Listings~\ref{lst:pncounter-monotonic}, \ref{lst:pncounter-progress}, and \ref{lst:pncounter-semantic}) establish core verification guarantees: monotonicity ensures information is never lost, progress preservation ensures updates propagate reliably, and semantic distinction formalizes that lattice order and computed value have fundamentally different properties.

The progress preservation property (Listing~\ref{lst:pncounter-progress}) proves that after a replica increments its positive component, merging with any other state will preserve that increment. This implements the formal requirement $\mathit{merge}(s, \mathit{apply}(s, v)) = \mathit{apply}(s, v)$ and ensures that any replica does not lose its local updates while synchronizing. 

\section{LWW Register: Deterministic Concurrent Write Resolution via Vector Clocks}\label{sec:lww}

LWW Register (Last-Write-Wins Register) addresses a fundamental problem in replicated state: when two replicas perform concurrent writes to the same register, how should the system deterministically select one value? The LWW Register uses vector clocks to establish a happens-before relation, enabling the system to order writes and apply a last-write-wins policy deterministically.

\subsection{The Challenge}

Consider replicas $R_1$ and $R_2$ concurrently writing different values to a register. The system must resolve these concurrent writes to a single value. Three approaches are possible:

\begin{itemize}
  \item \textbf{Arbitrary selection}: Choose one value randomly (no determinism).
  \item \textbf{Merge-based}: Combine both writes (not applicable for scalar values).
  \item \textbf{Last-write-wins}: Order writes by causality; apply whichever write came last (LWW Register).
\end{itemize}

LWW Register implements last-write-wins through vector clocks, which provide a total order on causally comparable events. When two writes are concurrent (not ordered by causality), the register applies a deterministic tie-breaking rule (lexicographic order on values). This requires the causal metadata provided by the Backbone framework (Chapter~\ref{chap:framework}).

\subsection{State and Operations}

\paragraph{State representation.} The LWW Register state consists of a value paired with its vector clock:
\[
\text{State} = (\text{value}, \text{clock}) \quad \text{where} \quad \text{value} \in \mathbb{V} \text{ and } \text{clock} \in \text{Clock}
\]
where $\mathbb{V}$ is the domain of values (typically identifiers or scalars).

The clock records the logical timestamp of the most recent write. An element is visible through the register lookup operation.

\paragraph{Key operations.} The LWW Register supports the following operations:
\begin{itemize}
  \item \textit{read}: $\mathit{read}((v, c)) = v$
  \item \textit{write}: $\mathit{write}((v_1, c_1), v_2, c_2) = (v_2, c_2)$
  \item \textit{merge}: $\mathit{merge}((v_1, c_1), (v_2, c_2)) = \text{if } c_1 \succ c_2 \text{ then } (v_1, c_1) \text{ else } (v_2, c_2)$
  \item \textit{tick}: $\mathit{tick}(c, r) = c[r \mapsto c(r)+1]$
\end{itemize}

Merge compares vector clocks: if one strictly dominates the other, that value is retained. If clocks are incomparable (concurrent writes), tie-breaking selects the lexicographically smaller value. The tick operation increments the local replica's clock before each write, maintaining causality.

\paragraph{Verified properties.} LWW Register's specification (188 lines in Athena) contains proofs of the following categories:

\begin{table}[h]
  \centering
  \small
  \begin{tabular}{lp{5cm}}
    \toprule
    \textbf{Category} & \textbf{Representative Theorems}\\
    \midrule
    Merge Structure & Commutativity, associativity, idempotence of merge \\
    \midrule
    Causality & Happens-before order ensures write precedence \\
    \midrule
    Convergence & Sequential writes guarantee convergence; concurrent writes use tie-breaking \\
    \bottomrule
  \end{tabular}
  \caption{Verified property categories for LWW Register}
  \label{tab:lww-theorems}
\end{table}

\paragraph{Clock and value preservation properties.} The LWW Register specification includes proven theorems for update and merge operations. We present two representative examples demonstrating how the register preserves vector clock structure and deterministic merge semantics.

First, Listing~\ref{lst:lww-update-value} shows the structure of the update operation:

\begin{lstlisting}[language=Athena,style=athena,label={lst:lww-update-value},caption={LWW Register: Update definition}]
declare lww-update: [LWWState Update ActorId] -> LWWState

assert* lww-update-def :=
  [((lww-update s (write v) a) =
    (pr v (clock-tick (second s) a)))]
\end{lstlisting}

This definition shows that updates pair the new value with a ticked clock.

Second, Listing~\ref{lst:lww-merge-clocks} proves that merge joins the vector clocks:

\begin{lstlisting}[language=Athena,style=athena,label={lst:lww-merge-clocks},caption={LWW Register: Merge joins clocks}]
conclude merge-clock-is-join :=
  (forall s1 s2 .
    (second (lww-merge s1 s2)) = (clock-join (second s1) (second s2)))
pick-any s1:LWWState s2:LWWState
  (!chain [(second (lww-merge s1 s2))
         = (second (pr (lww-select-value (first s1) (second s1) 
                                          (first s2) (second s2))
                       (clock-join (second s1) (second s2))))
           [lww-merge-def]
         = (clock-join (second s1) (second s2)) [Pr.pr-second]])
\end{lstlisting}

This establishes that merge computes the join (semilattice supremum) of the input clocks, which is the foundation for causal consistency.

These two theorems (Listings~\ref{lst:lww-update-value} and \ref{lst:lww-merge-clocks}) together ensure that LWW Register operations are well-formed: updates introduce new values with properly ticked clocks, and merges preserve the vector clock lattice structure that enables causal ordering.

\section{OR-Set: Composing CRDT State with Vector Clocks}\label{sec:orset}

The Observed-Remove Set (OR-Set) demonstrates a more complex interaction between CRDT semantics and causal ordering. Unlike state-based CRDTs (PN-Counter, G-Set), OR-Set requires vector clocks to implement add-wins semantics: when an add and remove operation occur concurrently, the add wins. This requires tracking causality to distinguish between ``remove followed by concurrent add'' (element present) versus ``add followed by causally-later remove'' (element absent).

\subsection{The Challenge: Re-addition Without Causality}

Consider two replicas performing operations on element $x$:
\begin{enumerate}
  \item $R_1$ adds $x$ at time $t_1$
  \item $R_2$ removes $x$ at time $t_2$ (concurrent with $t_1$)
  \item $R_1$ adds $x$ again at time $t_3$ (after seeing the removal)
\end{enumerate}

Without causality tracking, a pure state-based merge cannot distinguish the re-addition at $t_3$ from the original add at $t_1$. The OR-Set solves this by pairing each element with a unique tag generated from the replica's vector clock, creating a \emph{causal CRDT} that merges both set state and clock state.

\subsection{State Representation and Extended Merge}

OR-Set state pairs a set of tagged elements with a vector clock:
\[
\text{ExtendedORSetState} = (\text{ORSetState}, \text{Clock})
\]
where $\text{ORSetState}$ contains sets of adds (tagged elements) and removes (tags to tombstone).

The extended merge operation must merge \emph{both} the OR-Set state and the vector clocks, while maintaining the semilattice properties. This requires proving that the composition of two semilattices (set union semilattice and clock join semilattice) forms a valid semilattice.

\paragraph{Proof: Extended merge commutativity.} Listing~\ref{lst:orset-extended-merge} demonstrates how CRDT merge and vector clock join compose:\\

\begin{lstlisting}[language=Athena,style=athena,label={lst:orset-extended-merge},caption={OR-Set: Composing CRDT and vector clock commutativity}]
conclude extended-merge-commutative :=
  (forall s1 s2 .
    (orset-merge-with-context s1 s2) = (orset-merge-with-context s2 s1))
pick-any s1:ExtendedORSetState s2:ExtendedORSetState
    let {state1 := (first s1);
         clock1 := (second s1);
         state2 := (first s2);
         clock2 := (second s2);
         merged-state-12 := (orset-merge state1 state2);
         merged-state-21 := (orset-merge state2 state1);
         joined-clock-12 := (clock-join clock1 clock2);
         joined-clock-21 := (clock-join clock2 clock1);
         state-comm := (!chain [(orset-merge state1 state2)
                               = (orset-merge state2 state1)
                                   [Theory.merge-commutative-holds]]);
         clock-comm := (!chain [(clock-join clock1 clock2)
                               = (clock-join clock2 clock1)
                                   [VectorClocks.Properties.clock-join-commutative]]);
         result12 := (pr merged-state-12 joined-clock-12);
         result21 := (pr merged-state-21 joined-clock-21);
         pair-eq := (!chain [(pr merged-state-12 joined-clock-12)
                            = (pr merged-state-21 joined-clock-21)
                                [state-comm clock-comm]])}
      (!chain [(orset-merge-with-context s1 s2)
             = result12
             = result21  [pair-eq]
             = (orset-merge-with-context s2 s1)])
\end{lstlisting}

This proof combines the proven commutativity of OR-Set merge (from the CRDT specification) with the proven commutativity of vector clock join (from the framework's vector clock theory). The key insight is in the \texttt{pair-eq} step, which uses both \texttt{state-comm} and \texttt{clock-comm} to establish equality of the paired result.


\section{The Adapter Pattern in Practice}

Each CRDT is implemented by the adapter pattern shown in Figure~\ref{fig:adapter-pattern}. This section details how adapters modularize verification.

\subsection{Adapter Structure}

\begin{figure}[h]
\centering
\begin{tikzpicture}[
  box/.style={rectangle, draw, thick, minimum width=2.2cm, minimum height=1cm, align=center, font=\small},
  arrow/.style={->, thick, blue!70},
  stage/.style={rectangle, draw=none, minimum height=0.6cm, align=center, font=\tiny, text width=2cm}
]

  % Stage 1: Specification
  \node[box, fill=blue!20] (spec) at (0, 0) {\textbf{Specification}\\Module};
  \node[stage, below=0.3cm of spec] (spec-desc) {State type, ops,\\properties};
  
  % Arrow
  \draw[arrow] (1.3, 0) -- (2.2, 0);
  
  % Stage 2: Theory
  \node[box, fill=green!20] (theory) at (3.5, 0) {\textbf{Theory}\\Instantiation};
  \node[stage, below=0.3cm of theory] (theory-desc) {Proves axioms:\\ACI+Identity};
  
  % Arrow
  \draw[arrow] (4.8, 0) -- (5.7, 0);
  
  % Stage 3: Framework
  \node[box, fill=orange!20] (framework) at (7, 0) {\textbf{Framework}\\Instantiation};
  \node[stage, below=0.3cm of framework] (framework-desc) {PayloadOps\\implementation};
  
  % Arrow
  \draw[arrow] (8.3, 0) -- (9.2, 0);
  
  % Stage 4: Executable
  \node[box, fill=red!20] (exec) at (10.5, 0) {\textbf{Executable}\\Adapter};
  \node[stage, below=0.3cm of exec] (exec-desc) {Runnable\\instances};

\end{tikzpicture}
\caption{CRDT Adapter Template: The four-stage progression from formal specification to executable adapter. Each stage builds on previous verification: \emph{Specification module} defines state and operations; \emph{Theory instantiation} proves satisfaction of semilattice axioms (commutativity, associativity, idempotence, identity); \emph{Framework instantiation} provides \texttt{PayloadOps} implementations; \textit{Executable adapter} uses configuration machinery for testing and demonstration.}
\label{fig:adapter-pattern}
\end{figure}

\paragraph{Example: PN-Counter Adapter.}
The PN-Counter adapter (197 lines) consists of:

\begin{enumerate}
  \item \textbf{Specification module}: Defines state type, operations, and proves some theorems
  \item \textbf{Theory instantiation}: Proves ACI axioms for componentwise merge by inheriting GCounter proofs
  \item \textbf{Framework instantiation}: Registers with PayloadFramework, inheriting convergence guarantees
  \item \textbf{Demo module}: Instantiates replicas, performs operations, demonstrates convergence
\end{enumerate}

\subsection{Inheriting Properties via Adapters}

The theory/adapter pattern works by creating a renaming that maps CRDT operations to abstract semilattice operations, then automatically inheriting theorems. Listing~\ref{lst:gset-adapter} shows a concrete example from the GSet (Grow-only Set) implementation:

\begin{lstlisting}[language=Athena,style=athena,label={lst:gset-adapter},caption={GSet adapter: inheriting semilattice theorems via renaming}]
module GSetSpec {
  define-sort GSState := (Set.Set Ide)
  declare gs-merge: [GSState GSState] -> GSState
  declare gs-initial: GSState
  
  assert* gs-merge-is-union:=[((gs-merge ?s1 ?s2)=(Set.union ?s1 ?s2))]
  
  module Theory {
    conclude merge-commutative :=
      (forall s1 s2 . (gs-merge s1 s2) = (gs-merge s2 s1))
      pick-any s1:GSState s2:GSState
        (!chain [(gs-merge s1 s2)
               = (Set.union s1 s2) [gs-merge-is-union]
               = (Set.union s2 s1) [Set.union-commutes]
               = (gs-merge s2 s1) [gs-merge-is-union]])
    
    # Create adapter to abstract SemilatticeTheory
    define GSet-To-Semilattice :=
      (renaming |{SemilatticeTheory.join := gs-merge,
                  SemilatticeTheory.bot := gs-initial}|)
    
    # Inherit derived theorems automatically
    define gs-absorption :=
      (get-property SemilatticeTheory.absorption
                    GSet-To-Semilattice
                    SemilatticeTheory.theory)}}
\end{lstlisting}

Notice that after proving only the core ACI axioms and establishing the renaming, all other properties (absorption, monotonicity, convergence) are automatically available. The same \texttt{get-property} pattern applies to inherited theorems without requiring any manual proving.

\subsection{Executable Demonstration}

The adapter pattern is not merely a theoretical abstraction; it produces runnable code. Listing~\ref{lst:pncounter-demo} shows a demonstration script for PN-Counter that instantiates three replicas, performs concurrent increments and decrements, and verifies convergence:

\begin{lstlisting}[language=Athena,style=athena,label={lst:pncounter-demo},caption={PN-Counter executable demo: concurrent operations and convergence}]
load "../crdts/pncounter/pncounter_payload_adapter"
open PNCounterAdapter
open Payload

# Create initial configuration with 3 replicas
define actors := ('r1 :: 'r2 :: 'r3 :: nil)
define cfg0 := (pnc-init actors)

# Replica r1 increments by 3, broadcast and deliver
define cfg1 := (pnc-broadcast cfg0 ('r1 :: nil) 
                 (PNCounterSpec.inc-pos (S (S (S zero)))))
define cfg2 := (pnc-reduce cfg1)

# Replica r2 increments by 2, broadcast and deliver
define cfg3 := (pnc-broadcast cfg2 ('r2 :: nil) 
                 (PNCounterSpec.inc-pos (S (S zero))))
define cfg4 := (pnc-reduce cfg3)

# Replica r3 decrements by 1, broadcast and deliver
define cfg5 := (pnc-broadcast cfg4 ('r3 :: nil) 
                 (PNCounterSpec.inc-neg (S zero)))
define cfg6 := (pnc-reduce cfg5)

# Replica r1 decrements by 1, broadcast and deliver
define cfg7 := (pnc-broadcast cfg6 ('r1 :: nil) 
                 (PNCounterSpec.inc-neg (S zero)))
define cfg-final := (pnc-reduce cfg7)

# Evaluate final configuration
(eval (am cfg-final))
\end{lstlisting}

This demo script loads the verified PN-Counter specification. Each call to \texttt{pnc-broadcast} and \texttt{pnc-reduce} executes the merge operations whose properties were proven in Listings~\ref{lst:pncounter-monotonic}, \ref{lst:pncounter-progress}, and \ref{lst:pncounter-semantic}. The final \texttt{eval} displays the converged state across all replicas: each replica's state shows $(P, N) = (3, 1)$ with computed value $3 - 1 = 2$, demonstrating that monotonic state growth (all components increased) coexists with value semantics (the value may decrease). The framework thereby provides both machine-checked formal guarantees and executable instances for testing---demonstrating that verification need not practicality sacrifice.

\section{What We Verify: Summary}

The implementations chapter demonstrates that our verification system can prove:

\begin{itemize}
  \item \textbf{Structural properties}: Semilattice axioms (commutativity, associativity, idempotence) for merge operations
  
  \item \textbf{Monotonicity properties}: State components never decrease through merge operations (e.g., \texttt{positive-component-monotonic}, \texttt{negative-component-monotonic})
  
  \item \textbf{Progress properties}: Local updates are preserved through merges (e.g., \texttt{merge-preserves-pos-update})
  
  \item \textbf{Behavioral properties}: CRDT-specific invariants such as ``state grows but value may decrease'' (\texttt{state-grows-value-may-decrease}) and clock preservation in causal CRDTs (\texttt{merge-clock-is-join})
\end{itemize}

 Properties such as progress preservation are proven explicitly for PN-Counter (\texttt{merge-preserves-pos-update}, \texttt{merge-preserves-neg-update}), OR-Set (\texttt{merge-preserves-tag-sets}), G-Set (\texttt{merge-preserves-additions}), and 2P-Set (\texttt{remove-preserves-add-set}). For LWW-Register and G-Counter, progress follows from their proven merge properties via the theory/adapter pattern. In all scenarios the framework guarantees that updates applied at one replica remain visible after synchronization with concurrent replicas.

Each property is established through formal proof. All theorems presented in this chapter have been mechanically verified by Athena's proof engine.
